% ============================================================
%  Příprava na maturitu – Mechatronika
%  Kompiluj: pdflatex maturita_mechatronika.tex
%  Kódování: UTF-8
% ============================================================

\documentclass[a4paper,10pt]{article}

% --- Balíčky ---
\usepackage[T1]{fontenc}
\usepackage[utf8]{inputenc}
\usepackage[czech]{babel}
\usepackage[left=2.4cm, right=1.8cm, top=1.8cm, bottom=1.8cm]{geometry}
\usepackage{lmodern}
\usepackage{microtype}
\usepackage{parskip}
\usepackage{titlesec}
\usepackage{enumitem}
\usepackage{xcolor}
\usepackage{hyperref}
\usepackage{tabularx}
\usepackage{booktabs}
\usepackage{array}
\usepackage{tikz}
\usepackage{eso-pic}
\usepackage{mdframed}
\usepackage{amsmath}
\usepackage{amssymb}
\usepackage{pgfplots}
\pgfplotsset{compat=1.18}
\usetikzlibrary{arrows.meta, positioning, shapes.geometric, calc}

% --- Barvy ---
\definecolor{accent}{HTML}{03254E}
\definecolor{stripeblue}{HTML}{568EA3}
\definecolor{medgray}{HTML}{888888}
\definecolor{lightblue}{HTML}{EAF2F8}
\definecolor{lightyellow}{HTML}{FDFBE4}
\definecolor{lightgreen}{HTML}{EAF7EA}

% --- Modrý pruh vlevo ---
\AddToShipoutPictureBG{%
  \begin{tikzpicture}[remember picture, overlay]
    \fill[stripeblue]
      (current page.north west)
      rectangle
      ([xshift=10mm] current page.south west);
    \fill[accent!60]
      ([xshift=10mm] current page.north west)
      rectangle
      ([xshift=12mm] current page.south west);
  \end{tikzpicture}%
}

\hypersetup{colorlinks=true, urlcolor=accent, linkcolor=accent}

\titleformat{\section}
  {\large\bfseries\color{accent}}
  {\thesection.}
  {0.5em}
  {}
  [\color{accent}\titlerule]

\titlespacing{\section}{0pt}{16pt}{6pt}

\newmdenv[
  backgroundcolor=lightblue,
  linecolor=stripeblue,
  linewidth=1pt,
  roundcorner=4pt,
  innerleftmargin=8pt,
  innerrightmargin=8pt,
  innertopmargin=6pt,
  innerbottommargin=6pt,
  skipabove=4pt,
  skipbelow=4pt
]{pojmybox}

\newmdenv[
  backgroundcolor=lightyellow,
  linecolor=accent!40,
  linewidth=1pt,
  roundcorner=4pt,
  innerleftmargin=8pt,
  innerrightmargin=8pt,
  innertopmargin=6pt,
  innerbottommargin=6pt,
  skipabove=4pt,
  skipbelow=8pt
]{vysvetlenibox}

\newmdenv[
  backgroundcolor=lightgreen,
  linecolor=accent!30,
  linewidth=1pt,
  roundcorner=4pt,
  innerleftmargin=8pt,
  innerrightmargin=8pt,
  innertopmargin=6pt,
  innerbottommargin=6pt,
  skipabove=4pt,
  skipbelow=8pt
]{vzorcebox}

\pagestyle{plain}

% ============================================================
\begin{document}

% --- Záhlaví ---
\begin{minipage}[t]{0.75\textwidth}
    {\Huge\bfseries\color{accent} Příprava na maturitu}\\[4pt]
    {\large\color{stripeblue} Mechatronika}\\[2pt]
    \textcolor{medgray}{\small Suchánek Lukáš \quad SPŠ Prosek \quad 2026}
\end{minipage}
\hfill
\begin{minipage}[t]{0.22\textwidth}
    \raggedleft
    \begin{tikzpicture}
        \draw[color=stripeblue, rounded corners=4pt, line width=1pt]
            (0,0) rectangle (3,1.2)
            node[midway, align=center, color=accent, font=\small\bfseries]
            {Otázek: 22};
    \end{tikzpicture}
\end{minipage}

\vspace{6pt}
\noindent\color{accent}\rule{\linewidth}{1.5pt}
\vspace{2pt}

% ============================================================
\section{Akční členy -- elektrické, pneumatické a hydraulické mechanismy}

\begin{pojmybox}
\textbf{Klíčové pojmy:}
servomotor, krokový motor, DC motor, BLDC, solenoid,
přímočarý pohon, pneumatický válec, rozvaděč, hydraulický agregát,
proporcionální ventil, řízení výkonu
\end{pojmybox}

\begin{vysvetlenibox}
\textbf{Elektrické pohony -- srovnání}

\begin{center}
\begin{tabularx}{0.95\linewidth}{@{} l X X X @{}}
    \toprule
    \textbf{Typ motoru} & \textbf{Výhody} & \textbf{Nevýhody}
                        & \textbf{Aplikace} \\
    \midrule
    DC motor       & Jednoduché řízení (napětí = otáčky)
                   & Kartáče -- opotřebení & Modely, ventilátory \\
    BLDC motor     & Bez kartáčů, vysoká účinnost
                   & Složitější řízení (3-fázový) & Drony, e-kola \\
    Krokový motor  & Přesná poloha bez enkodéru
                   & Ztráta kroku při přetížení & 3D tiskárny, CNC \\
    Servomotor AC  & Velmi přesný, rychlý
                   & Drahý, složitý driver & CNC, robotika \\
    \bottomrule
\end{tabularx}
\end{center}

\begin{vzorcebox}
Moment DC motoru: $M = k_t \cdot I_a$,
otáčky: $n = \frac{U_a - I_a R_a}{k_e}$\\
Krokový motor -- krok: $\alpha = \frac{360°}{N_{kroků}}$,
typicky $\alpha = 1{,}8°$ (200 kroků/otáčku)\\
Výkon: $P = M \cdot \omega = M \cdot \frac{2\pi n}{60}$
\end{vzorcebox}

\textbf{Pneumatické systémy}\\
Pracují se stlačeným vzduchem (6--10\,bar).
Výhody: čisté, rychlé, nenákladné.
Nevýhody: hlučné, horší přesnost, energie se špatně zpět získává.

Základní prvky: kompresor $\to$ filtr-regulátor-maznice (FRL)
$\to$ rozvaděč (5/2, 5/3) $\to$ pneumatický válec.

\textbf{Hydraulické systémy}

\begin{vzorcebox}
Pascalův zákon: $p = \frac{F}{A}$ \quad [Pa = N/m²]\\
Síla hydraulického válce: $F = p \cdot A$\\
Příklad: $p = 200\,\text{bar}$, $A = 50\,\text{cm}^2$:
$F = 200 \times 10^5 \cdot 50 \times 10^{-4} = 100\,\text{kN}$
\end{vzorcebox}

\textbf{Srovnání pohonů}

\begin{center}
\begin{tabular}{@{} l c c c @{}}
    \toprule
    \textbf{Vlastnost} & \textbf{Elektrický}
                       & \textbf{Pneumatický} & \textbf{Hydraulický} \\
    \midrule
    Síla / hmotnost & Střední & Nízká & Velmi vysoká \\
    Přesnost polohy & Vysoká & Nízká & Střední \\
    Rychlost        & Vysoká & Velmi vysoká & Střední \\
    Účinnost        & 85--95\,\% & 20--40\,\% & 60--80\,\% \\
    Čistota         & Čistý & Čistý & Riziko úniku oleje \\
    \bottomrule
\end{tabular}
\end{center}
\end{vysvetlenibox}

% ============================================================

% ============================================================
\end{document}
